% l08_judgment.tex -- round 502.
% PRIME.RDAGGER.L08_JUDGMENT.01
% U1 BLOCKED(no-certifiable-gap@L>=0.45) /
% U2 LADDER_RESTS_REDIRECT_TO_BRIDGE_LANE /
% U3 ANTI_LIST_EXTENDED.
% Pure adjudication note.  No probe, no verifier, no new census, no pin.
% Every numeral is a float / dps-30 re-measurement made to ground the
% verdict; none of it is a sealed pin.
% NO RH CLAIM.  NO anti-RH CLAIM.
\documentclass[11pt]{article}

\usepackage[a4paper,margin=2.45cm]{geometry}
\usepackage{amsmath,amssymb,amsthm}
\usepackage{booktabs,array,microtype,xcolor}
\usepackage[T1]{fontenc}
\usepackage[colorlinks=true,linkcolor=blue!50!black,urlcolor=blue!50!black]{hyperref}

\newtheorem{lemma}{Lemma}
\newtheorem{proposition}{Proposition}
\newtheorem{contract}{Contract}
\theoremstyle{remark}
\newtheorem{remark}{Remark}
\newcommand{\QW}{Q_{\mathrm W}}
\newcommand{\blocked}{\textcolor{red!65!black}{BLOCKED}}
\setlength{\emergencystretch}{2em}

\title{\bfseries The $0.8$ rung: a judgment\\[2mm]
\large there is no gap to certify\\[2mm]
\large PRIME.RDAGGER.L08\_JUDGMENT.01 --- round 502}
\author{Stefan Hamann}
\date{August 31, 2026}

\begin{document}
\maketitle

\begin{abstract}\noindent
Adjudication of the round-496 no-go at $L=0.8$ and of its two named
obstructions (cut margin $-1.61\cdot10^{-4}$; noncompact prime
translations, off-block $0.9395$, $3\times$HS rest $27.367$).
The r496 numbers are first reproduced from an independent
implementation --- all four cutoff rows, $\kappa_w$,
$\|T_w\|_{\mathrm{HS}}=19.2177$, $r_{101}=9.12241$ --- which also
identifies the booking r496 actually used (prime translations retained
as operators).
\textbf{U1 \blocked{}$(\text{no-certifiable-gap}@L\ge0.45)$.}
Both obstructions are mis-specified, because the quantity they are
trying to make positive does not exist at the claimed size.  An
explicit odd test function of Legendre degree $\le13$ gives
\[
  \lambda_*(0.8)\ \le\ 1.8170\cdot10^{-14},
\]
verified in two independent representations: the $x$-space kernel form
at \texttt{dps=30} gives $1.81691694981\cdot10^{-14}$, and the
spectral side of the explicit formula, summed over the first $60$
zeros of $\zeta$ (all terms nonnegative), already gives
$1.442289883\cdot10^{-14}$, i.e.\ $79\%$ of it.
The r496 obstruction margins are therefore ten orders of magnitude
\emph{above} the target.  Measured ladder of ceilings:
$7.57\cdot10^{-3}$ $(L=0.30)$, $1.33\cdot10^{-3}$ $(0.34657)$,
$1.81\cdot10^{-4}$ $(0.40)$, $1.62\cdot10^{-5}$ $(0.45)$,
$9.34\cdot10^{-7}$ $(0.50)$, $1.61\cdot10^{-9}$ $(0.60)$,
$4.87\cdot10^{-13}$ $(0.70)$, below the float floor from $L=0.75$;
the archimedean floor meanwhile falls to $-0.92648$ at $L=0.8$ and
$-12.518$ at the r471 top rung $L_{16}=2.7726$, so the three channels
must cancel to $13.6$ digits at $L=0.8$ already.
Four candidate refinements are calibrated against these numbers and
all four die, each for a stated reason; the finest cutoff scan drives
the r496 sign deficit to $-7.44\cdot10^{-11}$ and its limit is the
ceiling itself.
\textbf{U2 LADDER\_RESTS / REDIRECT\_TO\_BRIDGE\_LANE:} the
$\lambda_*$ ladder is closed as a proof route, permanently, at every
rung beyond the prime-free zone.  The proved r493e assembly
\texttt{grid\_dense\_extension} consumes \emph{non-strict} positivity
uniform in the support length; a section-plus-remainder architecture
can never produce a non-strict inequality.
\textbf{U3} adds family F-VII (fixed-$L$ margin objects read as proof
targets) and six entries.
\textbf{Free diagnosis:} the compressed prime translations are exactly
$\tfrac12\times$ the adjacency operator of a path graph on
$\lceil2L/\log n\rceil$ vertices, so $\operatorname{spec}
(\operatorname{Re}R_{\log n})=\{\cos(\pi j/(m_n+1))\}$ --- measured
$\pm1/\sqrt2$, $\pm\tfrac12$, $\pm\tfrac12$ at $L=0.8$.  Their
obstruction is not noncompactness; it is incommensurability.
No RH claim, no anti-RH claim.
\end{abstract}

\medskip\noindent
\textbf{Status.}\enspace
\texttt{U1 BLOCKED(no-certifiable-gap@L>=0.45) /
U2 LADDER\_RESTS\_REDIRECT\_TO\_BRIDGE\_LANE /
U3 ANTI\_LIST\_EXTENDED}.
No promotion of an RH claim, no ledger row, no sealed pin.

\begin{center}
\fbox{\parbox{0.95\linewidth}{%
\textbf{Claim boundary.}
This is an adjudication note.  It contains no probe, no verifier, no
new census and no sealed pin; every numeral is a float or
\texttt{dps=30} re-measurement made to ground a verdict.
The inequality $\lambda_*(0.8)\le1.8170\cdot10^{-14}$ is an
\emph{upper} bound on a quantity that RH predicts to be
$\ge0$; a small positive $\lambda_*$ is exactly what RH predicts, and
an upper bound on it is \textbf{not} an anti-RH statement.  Both
representations return a \emph{positive} value for the witness and
they agree; nothing anomalous is observed.
The r494/r495 theorem $\lambda_*(0.3)\ge2.1\cdot10^{-3}$ and the
r471/r477 finite-subspace certificates are untouched.
No $L^*$ claim, no $R^\dagger$ claim, no RH claim, no anti-RH claim.}}
\end{center}

\section{Convention}

$\operatorname{supp}h\subset[-L,L]$, $g=h*\widetilde h$,
$g(0)=\|h\|_2^2$.  With $E$ the zero extension, $R_x=E^*\tau_xE$,
$a_n=\Lambda(n)/\sqrt n$ and $k(x)=e^{x/2}/\sinh x$, the r494 form is
\begin{equation}\label{eq:form}
  \QW(h)=\underbrace{\int_0^{2L}\!k(x)\bigl(g(0)-g(x)\bigr)dx-c_Lg(0)
          +\Pi(h)}_{\textstyle A(h)+\Pi(h)}
         \;-\;\underbrace{2\!\!\sum_{n\le e^{2L}}\!\!a_n\,g(\log n)}
              _{\textstyle P(h)} ,
\end{equation}
$\Pi(h)=2(C^2-S^2)$, $C=\langle h,\cosh(\cdot/2)\rangle$,
$S=\langle h,\sinh(\cdot/2)\rangle$, and
$\lambda_*(L)=\inf\{\QW(h)/\|h\|^2\}$ over the support class.
Two bookings of the prime term recur below:
\emph{(A)} $-P\succeq-2\sum a_n\cdot I$, dropping the positive
translation differences; \emph{(B)} $-P$ kept as the operator
$-2\sum a_n\operatorname{Re}R_{\log n}$.
The r469 scramble gate is live at this $L$ (the nodes carry the
literal $\log n$).  The r492 anti-list, including pads N1--N4, is in
force; in particular every Galerkin number below is typed as an
\textbf{upper} bound on an operator eigenvalue, which is the only
direction used.

\section{What was re-measured}\label{sec:remeas}

An independent implementation of \eqref{eq:form} was written for this
judgment (normalized Legendre Galerkin; the near-diagonal integrand
assembled through the cancellation-free split
$2(g(0)-g(x))=\int_{-L}^{L-x}|h(y)-h(y+x)|^2dy
+\int_{-L}^{-L+x}|h|^2+\int_{L-x}^{L}|h|^2$).

\subsection{Calibration against sealed numbers}

\begin{center}
\small
\begin{tabular}{@{}llll@{}}
\toprule
Quantity & this note & sealed & round \\
\midrule
$c_L(0.3)$ & $2.1924049111319992553$
  & $2.1924049111319992553$ & r495 \\
$\kappa_w(0.3)$, $[0.01,0.03]$ & $3.6980080402989774$
  & $3.6980080402989775$ & r494 \\
odd section, $L=0.3$ & $2.2256359052\cdot10^{-1}$
  & $2.2256359052\cdot10^{-1}$ & r494 \\
first Dirichlet, $L=0.3$ & $1.18025311091851\cdot10^{-2}$
  & $1.18025311091317\cdot10^{-2}$ & r494 \\
$c_L(0.8)$ & $3.5598735700355447$
  & $3.5598735700355446$ & r496 \\
$2\sum a_n$, $L=0.8$ & $2.9419735252236205$
  & $2.9419735252236205$ & r496 \\
constant mode, $L=0.8$ & $0.0802740549593$
  & $[0.0802740548742,0.0802740551870]$ & r477 \\
\bottomrule
\end{tabular}
\end{center}

\subsection{Reproduction of the r496 no-go}

\begin{center}
\small
\begin{tabular}{@{}lrrrr@{}}
\toprule
$[x_0,x_1]$ & $\kappa_w$
  & $\lambda_{\min}(P_{101}\QW^{(w)}P_{101})$
  & r496 & $\|T_w\|_{\mathrm{HS}}$ \\
\midrule
$[0.001,0.003]$ & $7.3770530633$ & $-1.610673\cdot10^{-4}$
  & $-1.6106733\cdot10^{-4}$ & $19.2177$ \\
$[0.002,0.006]$ & $6.6829061429$ & $-8.176622\cdot10^{-4}$
  & $-8.1766226\cdot10^{-4}$ & $13.5946$ \\
$[0.003,0.010]$ & $6.1992956339$ & $-2.284102\cdot10^{-3}$
  & $-2.2841017\cdot10^{-3}$ & $10.6441$ \\
$[0.010,0.030]$ & $5.0654766992$ & $-3.032258\cdot10^{-2}$
  & $-3.0322577\cdot10^{-2}$ & $6.0742$ \\
\bottomrule
\end{tabular}
\end{center}

All four rows, all four $\kappa_w$, and
$r_{101}=9.12241$ against the sealed $9.122410180$ (hence
$3r_{101}=27.3672$ against $27.367230540$) reproduce.
\textbf{Booking identified:} the reproduction only matches under
booking (B).  Under booking (A) the same sections give
$-3.8688$, four orders of magnitude lower.  This matters: r496's
surviving margin is the one that keeps the prime translations
\emph{in} the error block, which is exactly why its G3 then failed.

\section{U1 --- the ceiling}\label{sec:ceiling}

\subsection{The witness}

The uncut Legendre Galerkin $\lambda_{\min}$ of \eqref{eq:form} at
$L=0.8$ is $6.07\cdot10^{-15}$ at $n=101$ and does not move with $n$
($7.79\cdot10^{-14}$, $8.73\cdot10^{-15}$, $1.10\cdot10^{-14}$,
$7.26\cdot10^{-15}$ at $n=25,51,101,151$): the value is at the
float floor, not converging to anything resolvable.
Its eigenvector is odd to $1.5\cdot10^{-8}$ of its mass, with
dominant normalized Legendre coefficients
\[
 (c_1,c_3,c_5,c_7,c_9,c_{11},c_{13})
 \approx(0.4099,\,0.6641,\,0.5495,\,0.2847,\,0.0891,\,0.0103,\,0.00456),
\]
and $99.9999\%$ of its mass below degree $20$.  Evaluated as an
explicit function --- one polynomial, no Galerkin, no eigenvalue ---
in the $x$-space form at \texttt{dps=30}:
\[
 \boxed{\ \frac{\QW(h)}{\|h\|_2^2}=1.81691694981\cdot10^{-14}\ }
\]
against channel values $A(h)=-0.24300$, $\Pi(h)=-0.02894$,
$-P(h)=+0.27194$: a $13.6$-digit cancellation.

\subsection{The independent confirmation}

Under \eqref{eq:form} the program's own normalization makes
$\QW(h)=\sum_\rho\widehat g(\rho)$, which on the critical line is
$\sum_\gamma|\widehat h(\gamma)|^2$ --- a sum of \emph{nonnegative}
terms.  Every partial sum is therefore a hard lower bound, and it is
a third representation, independent of both the $x$-space and the
frequency form.  With the first $60$ zeros of $\zeta$
($\gamma_{60}=163.0307$):

\begin{center}
\small
\begin{tabular}{@{}lrrr@{}}
\toprule
Test function & $x$-space $\QW/\|h\|^2$
  & spectral, $60$ zeros & captured \\
\midrule
$L=0.8$ constant (r477 control) & $0.0802740549593$
  & $0.07050676638$ & $88\%$ \\
$L=0.3$ even minimiser (r494 control) & $0.0075718518899$
  & $0.007040224407$ & $93\%$ \\
$L=0.6$ witness & $1.62318271558\cdot10^{-9}$
  & $1.39688596\cdot10^{-9}$ & $86\%$ \\
$L=0.8$ witness & $1.81691694981\cdot10^{-14}$
  & $1.442289883\cdot10^{-14}$ & $79\%$ \\
\bottomrule
\end{tabular}
\end{center}

The two controls behave as they must (monotone from below, under the
total).  The $L=0.8$ witness is bracketed:
$1.4423\cdot10^{-14}\le\QW(h)/\|h\|^2=1.8169\cdot10^{-14}$.
Term by term the low zeros are minuscule --- the $\gamma_1$ term is
$1.06\cdot10^{-26}$, i.e.\ $|\widehat h(\gamma_1)|\approx
7.3\cdot10^{-14}$ against $\|\widehat h\|_2^2=2\pi$.
The near-null direction is real, not a numerical artefact, and it is
a genuine near-annihilator of the low zeros.

\begin{proposition}[the ceiling]\label{prop:ceiling}
$\lambda_*(0.8)\le1.8170\cdot10^{-14}$, witnessed by one explicit
odd polynomial.  Sealing it needs one interval quadrature of that
polynomial; it is not sealed here and the verdict does not need it.
\end{proposition}

\subsection{The ladder}

Galerkin ceilings ($n=81$; upper bounds on $\lambda_*$, which is the
correct logical type) against the archimedean-only floor
$\lambda_{\min}(A+\Pi)$, and the resulting cancellation depth at the
witness:

\begin{center}
\small
\begin{tabular}{@{}rrrrr@{}}
\toprule
$L$ & nodes & $\lambda_{\min}(A+\Pi)$
  & $\lambda_*(L)\le$ & digits \\
\midrule
$0.30000$ & $0$ & $+7.5715\cdot10^{-3}$ & $7.5715\cdot10^{-3}$ & $2.1$ \\
$0.34657$ & $0$ & $+1.3294\cdot10^{-3}$ & $1.3294\cdot10^{-3}$ & $2.9$ \\
$0.40000$ & $1$ & $-7.8049\cdot10^{-2}$ & $1.8136\cdot10^{-4}$ & $3.8$ \\
$0.45000$ & $1$ & $-2.0521\cdot10^{-1}$ & $1.6166\cdot10^{-5}$ & $4.9$ \\
$0.50000$ & $1$ & $-3.2224\cdot10^{-1}$ & $9.3448\cdot10^{-7}$ & $6.1$ \\
$0.60000$ & $2$ & $-5.3588\cdot10^{-1}$ & $1.6096\cdot10^{-9}$ & $8.9$ \\
$0.70000$ & $3$ & $-7.3394\cdot10^{-1}$ & $4.8745\cdot10^{-13}$ & $12.5$ \\
$0.80000$ & $3$ & $-9.2648\cdot10^{-1}$ & float floor & $13.6$ \\
$1.00000$ & $5$ & $-1.3241$ & float floor & --- \\
$1.50000$ & $12$ & $-2.6504$ & float floor & --- \\
$2.77260$ & $70$ & $-12.5178$ & float floor & --- \\
\bottomrule
\end{tabular}
\end{center}

Three readings, all load-bearing for the verdict.

\begin{enumerate}
\item \textbf{The collapse is $\approx3$ decades per $\Delta L=0.1$}
and it starts exactly where the prime nodes turn on.  By $L=0.45$ the
required relative precision is already below $10^{-5}$; by $L=0.7$ it
is $10^{-12.5}$.
\item \textbf{The archimedean channel goes negative at $L\approx0.37$}
and reaches $-0.92648$ at $L=0.8$.  Beyond the prime-free zone the
prime term is not a perturbation to be bounded --- it \emph{carries}
the positivity, and it must do so to the last digit.
\item \textbf{r494 succeeded because $L=0.3$ needs $2.1$ digits.}
That is the whole difference between r494 and r496, and it was
invisible in the r492 ranking, which listed ``same machine up the
r471 schedule, $2$--$5$ rounds''.  That ranking entry is hereby
withdrawn.
\end{enumerate}

\subsection{The candidate table}

Each candidate was calibrated against the numbers above
\emph{before} being named or killed, per the r492 standard.

\begin{center}
\small
\begin{tabular}{@{}llp{0.54\linewidth}@{}}
\toprule
Cand. & Verdict & Ground \\
\midrule
(K1) finer cutoff window
  & \textsc{kill}, transposed
  & the deficit really does vanish --- and its limit is the ceiling
    (\S\ref{sec:k1}) \\
(K2) prime-free error block via a scalar
  & \textsc{kill} $\times2$
  & $\lambda_{\min}(A+\Pi)=-0.92648<0$ kills every scalar or
    potential booking (\S\ref{sec:k2}) \\
(K3) primes as an exact finite-parametric object
  & \textsc{kill}, structural
  & right per node, impossible jointly: $\log3/\log2\notin\mathbb Q$
    (\S\ref{sec:k3}) \\
(K4) strip-booked Loewner cutoff
  & \textsc{kill} at the witness
  & the witness has $h(\pm L)^2=9.84\cdot10^{-14}$; the strip buys
    nothing (\S\ref{sec:k4}) \\
\bottomrule
\end{tabular}
\end{center}

\subsection{(K1) The cutoff scan}\label{sec:k1}

The r496 window grid had four rows.  Continuing it:

\begin{center}
\small
\begin{tabular}{@{}lrrr@{}}
\toprule
$[x_0,x_1]$ & $\kappa_w$
  & $\lambda_{\min}(P_{101}\QW^{(w)}P_{101})$
  & $r_{101}$ \\
\midrule
$[10^{-3},3\cdot10^{-3}]$ & $7.3771$ & $-1.6107\cdot10^{-4}$ & $9.12$ \\
$[3\cdot10^{-4},10^{-3}]$ & $8.5048$ & $-1.3454\cdot10^{-5}$ & $21.02$ \\
$[10^{-4},3\cdot10^{-4}]$ & $9.6805$ & $-9.0685\cdot10^{-7}$ & $48.16$ \\
$[3\cdot10^{-5},10^{-4}]$ & $10.8077$ & $-9.2299\cdot10^{-8}$ & $94.96$ \\
$[10^{-5},3\cdot10^{-5}]$ & $11.9832$ & $-6.9242\cdot10^{-9}$ & $182.65$ \\
$[10^{-6},3\cdot10^{-6}]$ & $14.2858$ & $-7.4437\cdot10^{-11}$ & $601.58$ \\
\bottomrule
\end{tabular}
\end{center}

So O1 is not a wall: the sign deficit falls monotonically and would
cross zero.  It crosses at the uncut value --- the ceiling of
Proposition~\ref{prop:ceiling}.  The scan is converging to the right
answer and the right answer is $10^{-14}$.
Quantitatively, at the witness the booked cutoff loss is dominated by
its Dirichlet part, $\tfrac14x_0^2\|h'\|^2$ with
$\|h'\|^2=34.5442$, i.e.\ $8.636\,x_0^2$.  For a certificate to
remain positive \emph{at the witness} it must satisfy
$8.636\,x_0^2<1.8170\cdot10^{-14}$, hence
\[
 x_0<4.59\cdot10^{-8}
 \quad\Longrightarrow\quad
 n\ \gtrsim\ 2L/x_0=3.5\cdot10^{7},
 \qquad \|T_w\|_{\mathrm{HS}}\approx2.8\cdot10^{3},
\]
and $3r_n<1.8\cdot10^{-14}$ then demands a relative Hilbert--Schmidt
tail of $2\cdot10^{-18}$, pushing $n$ higher again.  Against r494's
$401^2=1.6\cdot10^5$ interval-certified entries this is
$\gtrsim10^{15}$ of them.  \textbf{Kill, transposed:} the mechanism
is sound, the target is not there.

\subsection{(K2) A prime-free error block}\label{sec:k2}

This is O2's own question --- can \eqref{eq:form} be split into an
archimedean part certifiable the $0.3$ way, plus prime Loewner
positivity, minus a prime scalar, with the error block prime-free?
Two sharpenings of the scalar were found and measured, and both are
irrelevant for the same reason.

\emph{The almost-periodic sup.}  The crude scalar $2\sum a_n=2.9420$
is not the sharp one.  Since compression never increases a
multiplier's sup,
\[
 -P\ \succeq\ -\pi_L\,I,\qquad
 \pi_L=\sup_{t}\Bigl(-2\!\!\sum_{n\le e^{2L}}\!\!a_n\cos(t\log n)\Bigr).
\]
At $L=0.8$ this is available in closed form, because $\log4=2\log2$:
with $c=\cos(t\log2)$ the $n\in\{2,4\}$ part is the quadratic
$a_2c+a_4(2c^2-1)$, minimal at $c=-a_2/(4a_4)=-1/(2\sqrt2)$, while
$\cos(t\log3)=-1$ is attainable in the closure because
$\log3/\log2\notin\mathbb Q$ (Weyl).  Hence
$\pi_{0.8}=2.13500217689506$, confirmed by a brute-force scan
($2.134999$ over $t\le4000$): a gain of $0.8070$ over $2\sum a_n$.

\emph{The boundary-strip minorant.}  Better still, from the exact
zero-extension identity,
$\|h-\tau_ah\|^2\ge\int_{-L}^{-L+a}|h|^2+\int_{L-a}^{L}|h|^2$, so
\[
 -P\ \succeq\ -\Bigl(\sum_n a_n\Bigr)I+M_V,\qquad
 \sum_n a_n=1.4709867626,
\]
with $M_V$ multiplication by an explicit nonnegative step function,
levels $(0.4907286,\,0.9808577,\,0.3465736,\,0)$ from the centre
outwards.  This beats $\pi_L$ and is exact.

\emph{Why both are dead.}  Any such split bounds the certificate by
$\lambda_{\min}(A+\Pi)-(\text{prime scalar})+\min V$.  Measured,
$\lambda_{\min}(A+\Pi)=-0.9264812$ at $L=0.8$ (converged:
$-0.9264551$, $-0.9264781$, $-0.9264812$, $-0.9264823$ at
$n=25,51,101,151$), $\min V=0$ near the boundary, and
$\sup V=0.9809<\sum a_n$.  Every scalar or potential booking of the
prime term therefore returns a certificate between $-2.40$ and
$-3.87$; the crude end of that range is exactly the measured
booking-(A) section $-3.8688$.
There is a second, independent kill: even a perfect such split would
have to land below $1.82\cdot10^{-14}$.
\textbf{Kill $\times2$.}  The answer to O2's decomposition question is
a measured \emph{no}: beyond the prime-free zone the error block
cannot be made prime-free, because the primes are not an error.

\subsection{(K3) The primes as a finite-parametric object}\label{sec:k3}

O2's second intuition --- three nodes, rank-like structure in the
right coordinates --- is correct per node, and is worth recording
exactly.

\begin{lemma}[compressed translations are path graphs]\label{lem:path}
Let $a>0$, $m=\lceil 2L/a\rceil$.  Then $R_a$ is a partial isometry
with $R_a^{\,j}=R_{ja}$ and $R_a^{\,m}=0$, and
$\operatorname{Re}R_a$ is unitarily $\tfrac12$ times the adjacency
operator of the path graph on $m$ vertices, tensored with the
multiplicity of the cell shift.  Hence
$\operatorname{spec}(\operatorname{Re}R_a)
 \subset\{\cos(\pi j/(m+1))\}_{j=1}^{m}\cup\{0\}$.
\end{lemma}

Measured at $L=0.8$, $n=101$: $\operatorname{Re}R_{\log2}$
($m=3$) has $\lambda_{\max}=0.707107=\cos(\pi/4)$,
$\operatorname{Re}R_{\log3}$ and $\operatorname{Re}R_{\log4}$
($m=2$) have $\lambda_{\max}=0.500000=\cos(\pi/3)$, each to all
printed digits.  So each prime block is bounded, noncompact, and
\emph{exactly diagonalizable in closed form}: r496's off-block
$0.9395$ is not a smallness failure, and calling these terms
``noncompact'' understates what is true of them individually.

The obstruction is joint.  An exact finite reduction handling all
nodes at once needs one cell length dividing every $\log n$; for
$n=2,3$ that requires $\log3/\log2\in\mathbb Q$.  It is not.  What is
left is the finite-section theory of an almost-periodic symbol on an
interval --- precisely the object whose frequency-side incarnation
r492 closed permanently.  \textbf{Kill, structural}, and this is the
correct diagnosis of the r496 measurement.

\subsection{(K4) The strip-booked cutoff}\label{sec:k4}

The refined Loewner booking O1 asks for exists:
\[
 \int_0^{x_0}\!k(x)\bigl(I-\operatorname{Re}R_x\bigr)dx
 \ \succeq\ \tfrac12M_{W},\qquad
 W(y)=\int_{d(y)}^{x_0}\!\!k,\quad d(y)=\operatorname{dist}(y,\{\pm L\}),
\]
worth $\approx x_0\,h(\pm L)^2$, which for a generic mode at
$x_0=10^{-3}$ is $\approx6\cdot10^{-4}$ and would indeed dominate the
$-1.61\cdot10^{-4}$ deficit.  Two independent reasons it is worthless
here.  First, mechanically: $M_W$ is an unbounded positive
multiplication operator, and the block bound cannot digest it
($\|P M_WQ\|$ is not small); the best finite-rank minorant supported
on the two inner half-strips carries only $\int_{x_0/2}^{x_0}k=\log2$
and lands at $2.2\cdot10^{-4}$, marginal even against the wrong
target.  Second, and decisively: \emph{the witness has no endpoint
mass}.  Measured, $h(L)=3.14\cdot10^{-7}$, i.e.\
$h(\pm L)^2=9.84\cdot10^{-14}$ --- the same order as the ceiling
itself.  The near-null direction essentially vanishes at the
endpoints, so the entire strip booking is worth $10^{-13}$ there.
\textbf{Kill at the witness.}

\subsection{The verdict}

\begin{proposition}[round-502 judgment]\label{prop:blocked}
Beyond the prime-free zone the r492/r494 architecture --- booked
Loewner cutoff, Legendre finite section, remainder charged at
$r+2r$ --- is \blocked{}.  It certifies statements of the form
$\lambda_*(L)\ge c>0$ by subtracting three strictly positive booked
quantities from a scalar balance.  At $L=0.8$ the target is
$\le1.8170\cdot10^{-14}$, the finest booking measured anywhere in the
$(x_0,n)$ plane is $7.44\cdot10^{-11}$, and forcing the cutoff loss
alone under the target requires $n\gtrsim3.5\cdot10^{7}$ modes.
From $L\approx0.45$ upwards the required relative precision exceeds
five digits and grows by $\approx3$ decades per $\Delta L=0.1$; from
$L\approx0.75$ the target is below double precision.  No refinement of
the booking changes this, because the deficit is not in the booking.
\end{proposition}

\section{U2 --- the resource decree}

\subsection{The ladder was asking the wrong question}

The decisive observation is not numerical.  The r493e Lean assembly
\begin{center}
\texttt{grid\_dense\_extension} :
$(\forall f.\ 0\le\texttt{weilForm}\ f)\Rightarrow
 \forall F\ \text{admissible}.\ 0\le\texttt{fullWeilForm}\ F$
\end{center}
is proved sorry-free, and its hypothesis is \textbf{non-strict}
positivity on the grid class, uniform in the support length.  That is
what the assembled proof consumes.  A fixed-$L$ gap $\lambda_*(L)\ge
c>0$ is therefore both \emph{stronger} than what the assembly needs
at that $L$ and \emph{weaker} than what RH needs, since RH needs all
$L$.  And a section-plus-remainder architecture can never deliver the
non-strict statement: every remainder charge it makes is strictly
positive, so it can only ever prove a strict inequality --- exactly
the kind that \S\ref{sec:ceiling} shows does not survive.
The $\lambda_*$ ladder is thus not merely hard beyond $L=0.3465$; it
is aimed off-target.

\subsection{Decree}

\begin{enumerate}
\item \textbf{The $\lambda_*$ ladder is closed as a proof route},
permanently, for every $L>(\log2)/2$.  No round may reopen a booked
cutoff, a Legendre finite section with an HS remainder charge, or any
variant of the r494 three-gate contract at a rung beyond the
prime-free zone.  Its legacy is the ceiling curve of
\S\ref{sec:ceiling}, which is a genuine and reusable measurement.
\item \textbf{r494/r495 stand.} $\lambda_*(0.3)\ge2.1\cdot10^{-3}$ is
untouched, doubly checked, and now also \emph{explained}: $L=0.3$ is
the last rung needing only $2.1$ digits.
\item \textbf{Resources go to the bridge lane} (r497--r501: [2c]
spectral-sum analysis, then [2d] the contour form).  It is the long
pole, it is producing at a steady cadence, and it is the lane whose
output the assembly actually consumes.  Nothing in this note touches
it.
\item \textbf{Any future internal analysis round} must target the
non-strict, $L$-uniform grid-class statement
$0\le\texttt{weilForm}\,f$ by an exact or structural route --- a sum
of squares, a factorization, an identity --- never a numeric margin.
\item \textbf{Feasibility pre-gate, mandatory.} Before any scheme is
proposed at a support length $L$, the cancellation depth at $L$ must
be measured first (the ``digits'' column of \S\ref{sec:ceiling}).
A scheme whose booked losses exceed $10^{-\text{digits}}$ is refused
without further review.  This gate would have refused r496 before it
ran.
\item \textbf{F2b} (r471 density plus continuity modulus) is
superseded for the fixed-support direction by the proved r493e
density bridge, and is not revived by this note.
\item The r469 scramble gate stays in force; the measurements here are
node-literal and scramble-sensitive.
\end{enumerate}

\section{U3 --- anti-list}\label{sec:u3}

Families F-I--F-VI and pads N1--N4 of r492 stay in force.  N1 is not
violated here: Galerkin numbers are used only as upper bounds on
$\lambda_*$, which is their correct logical type, and the
load-bearing ceiling is a single high-precision Rayleigh quotient
bracketed from below by a nonnegative spectral partial sum.
New family:

\begin{center}
\small
\begin{tabular}{@{}llp{0.50\linewidth}@{}}
\toprule
Family & Rounds & Members \\
\midrule
F-VII fixed-$L$ margin objects read as
  proof targets
  & r471, r477, r492, r496, \textbf{r502}
  & see N5--N10 \\
\bottomrule
\end{tabular}
\end{center}

\begin{center}
\small
\begin{tabular}{@{}lp{0.86\linewidth}@{}}
\toprule
N5 & a fixed-$L$ finite-subspace value read as an approximation to
     $\lambda_*(L)$.  Measured: the r471 $L=0.8$ hint
     $2.645\cdot10^{-4}$ and $\mu_{\rm lo}=1.3225\cdot10^{-4}$ exceed
     $\lambda_*(0.8)$ by more than ten orders of magnitude.  They
     remain valid as subclass statements; r496 said so, and this is
     the number \\
N6 & any cutoff, window or booking margin quoted at $L>0.4$ without
     the measured $\lambda_*$ ceiling beside it.  Type specimen: the
     r496 reading of $-1.61\cdot10^{-4}$ as ``agonizingly close'' ---
     the sign deficit is real and it is $10^{10}$ times the target \\
N7 & the r492 U2-ranking entry ``same machine up the r471 schedule,
     $2$--$5$ rounds, yield $\lambda_*(L)\ge c$ at $L=0.8$ and
     beyond'' --- refuted and withdrawn \\
N8 & the prime term booked as a scalar ($2\sum a_n$), as an
     almost-periodic sup ($\pi_L$), or as scalar-plus-potential
     ($\sum a_n$ with $M_V$).  All three give negative certificates,
     because $\lambda_{\min}(A+\Pi)(0.8)=-0.92648<0$ \\
N9 & the compressed prime translations treated as a joint finite-rank
     or finite-section object.  Per node they are exactly
     $\tfrac12\,\mathrm{adj}(P_{m_n})$ (Lemma~\ref{lem:path});
     jointly they are incommensurable, $\log3/\log2\notin\mathbb Q$ \\
N10 & a strict-gap architecture applied to a non-strict target.
     Every remainder charge is positive, so such an architecture
     cannot prove $\QW\succeq0$ without a gap --- and beyond
     $L\approx0.75$ no gap survives double precision \\
\bottomrule
\end{tabular}
\end{center}

\section{Free diagnoses}

\begin{enumerate}
\item \textbf{Lemma~\ref{lem:path}} (exact spectrum of the compressed
translation) is new here and is reusable wherever a truncated shift
appears; it replaces ``bounded but noncompact'' with a closed form.
\item \textbf{$\lambda_*(L)$ is the frame lower bound of the
$\zeta$-zero set for $PW_L$.}  Under \eqref{eq:form},
$\lambda_*(L)=\inf\sum_\gamma|\widehat h(\gamma)|^2$ over unit-norm
type-$L$ functions.  The ceiling curve of \S\ref{sec:ceiling} is the
program's first quantitative picture of that frame constant, and it
explains the collapse: once the type is large enough, the low zeros
can be near-annihilated ($|\widehat h(\gamma_1)|\approx7\cdot10^{-14}$
at $L=0.8$).  Nothing about this is evidence for or against RH; it is
the reason a gap-based route cannot climb.
\item \textbf{Cross-representation health.}  The $x$-space form, the
frequency form (r496 G1, $5.5\cdot10^{-14}$) and now the spectral form
agree on every test function tried, including two controls pinned in
earlier rounds.  The witness value is positive in both
representations that can produce a number.  No anomaly.
\item \textbf{The r496 booking.}  Reproduction shows r496's surviving
row is booking (B).  Its G2/G3 tension was therefore structural from
the start: (B) is the only booking with a chance of positivity, and
(B) is the only booking that puts noncompact translations in the
error block.  The two obstructions O1 and O2 were the same
obstruction seen twice.
\end{enumerate}

\section{What is not claimed}

No $\lambda_*(L)\ge c$ for any $L>0.3$.  No claim that
$\lambda_*(0.8)=0$: the measurement is an upper bound
$1.8170\cdot10^{-14}$ and both representations return a positive
value.  No counterexample, no anti-positivity statement, and in
particular nothing against RH --- a small positive frame constant is
what RH predicts.  No Galerkin value is promoted to an operator
value.  The r471/r477 subclass certificates and the r494/r495
fixed-support theorem stand unchanged.  Every numeral here is a float
or \texttt{dps=30} re-measurement, not a sealed pin.
No $L^*$ claim, no $R^\dagger$ claim, no RH claim, no anti-RH claim.

\end{document}
